% Chunk: §3.4 Rates and Cross-Sections
% Source: Weinberg QFT Vol I, printed pp. 134-140 (PDF pp. 157-163)
% Status: transcribed and source-reviewed
Similarly, the observation in 1964 of small violations of CP conservation in the weak
interactions together with the assumed invariance of all interactions under CPT allowed the
immediate inference that the weak interactions also do not exactly conserve T. This has since
been verified~\hyperref[ref:3.18]{[18]} by more detailed studies of the $K^0$--$\bar K^0$ system, but it has so far
been impossible to find other direct evidence of the failure of invariance under time-reversal.

\subsection{Rates and Cross-Sections}

The $S$-matrix $S_{\beta\alpha}$ is the probability amplitude for the transition
$\alpha\to\beta$, but what does this have to do with the transition rates and cross-sections
measured by experimentalists? In particular, \eqref{eq:3.3.2} shows that $S_{\beta\alpha}$ has
a factor $\delta^4(p_\beta-p_\alpha)$, which ensures the conservation of the total energy
and momentum, so what are we to make of the factor $[\delta^4(p_\beta-p_\alpha)]^2$
in the transition probability $|S_{\beta\alpha}|^2$? The proper way to approach these
problems is by studying the way that experiments are actually done, using wave packets
to represent particles localized far from each other before a collision, and then following
the time-history of these superpositions of multi-particle states. In what follows we will
instead give a quick and easy derivation of the main results, actually more a mnemonic than
a derivation, with the excuse that (as far as I know) no interesting open problems in physics
hinge on getting the fine points right regarding these matters.

We consider our whole system of physical particles to be enclosed in a large box with a
macroscopic volume $V$. For instance, we can take this box as a cube, but with points on
opposite sides identified, so that the single-valuedness of the spatial wave function requires
the momenta to be quantized
\begin{equation}
 \mathbf{p}=\frac{2\pi}{L}(n_1,n_2,n_3),
 \tag{3.4.1}\label{eq:3.4.1}
\end{equation}
where the $n_i$ are integers, and $L^3=V$. Then all three-dimensional delta-functions become
\begin{equation}
 \delta_V^3(\mathbf{p}'-\mathbf{p})\equiv
 \frac{1}{(2\pi)^3}\int_V d^3x\,e^{i(\mathbf{p}-\mathbf{p}')\cdot\mathbf{x}}
 =\frac{V}{(2\pi)^3}\delta_{\mathbf{p}',\mathbf{p}},
 \tag{3.4.2}\label{eq:3.4.2}
\end{equation}
where $\delta_{\mathbf{p}',\mathbf{p}}$ is an ordinary Kronecker delta symbol, equal to one if the
subscripts are equal and zero otherwise. The normalization condition \eqref{eq:3.1.2} thus implies
that the states we have been using have scalar products in a box which are not just sums of
products of Kronecker deltas, but also contain a factor $[V/(2\pi)^3]^N$, where $N$ is the
number of particles in the state. To calculate transition probabilities we should use states
of unit norm, so let us introduce states normalized approximately for our box
\begin{equation}
 \ket{\alpha}^{\mathrm{Box}}\equiv
 \left[(2\pi)^3/V\right]^{N_\alpha/2}\ket{\alpha},
 \tag{3.4.3}\label{eq:3.4.3}
\end{equation}
with norm
\begin{equation}
 {}^{\mathrm{Box}}\!\braket{\beta}{\alpha}^{\mathrm{Box}}=\delta_{\beta\alpha},
 \tag{3.4.4}\label{eq:3.4.4}
\end{equation}
where $\delta_{\beta\alpha}$ is a product of Kronecker deltas, one for each three-momentum,
spin, and species label, plus terms with particles permuted. Correspondingly, the $S$-matrix
may be written
\begin{equation}
 S_{\beta\alpha}=\left[V/(2\pi)^3\right]^{(N_\beta+N_\alpha)/2}
 S^{\mathrm{Box}}_{\beta\alpha},
 \tag{3.4.5}\label{eq:3.4.5}
\end{equation}
where $S^{\mathrm{Box}}_{\beta\alpha}$ is calculated using the states \eqref{eq:3.4.3}.

Of course, if we just leave our particles in the box forever, then every possible transition
will occur again and again. To calculate a meaningful transition probability we also have to
put our system in a `time box'. We suppose that the interaction is turned on for only a time
$T$. One immediate consequence is that the energy-conservation delta function is replaced
with
\begin{equation}
 \delta_T(E_\alpha-E_\beta)=\frac{1}{2\pi}\int_{-T/2}^{T/2}
 \exp\bigl(i(E_\alpha-E_\beta)t\bigr)\,dt.
 \tag{3.4.6}\label{eq:3.4.6}
\end{equation}
The probability that a multi-particle system, which is in a state $\alpha$ before the
interaction is turned on, is found in a state $\beta$ after the interaction is turned off, is
\begin{equation}
 P(\alpha\to\beta)=|S^{\mathrm{Box}}_{\beta\alpha}|^2
 =\left[(2\pi)^3/V\right]^{N_\beta+N_\alpha}|S_{\beta\alpha}|^2.
 \tag{3.4.7}\label{eq:3.4.7}
\end{equation}
This is the probability for a transition into one specific box state $\beta$. The number of
one-particle box states in a momentum-space volume $d^3p$ is $Vd^3p/(2\pi)^3$, because this
is the number of triplets of integers $n_1,n_2,n_3$ for which the momentum \eqref{eq:3.4.1} lies in the
momentum-space volume $d^3p$ around $\mathbf{p}$. We shall define the final-state interval
$d\beta$ as a product of $d^3p$ for each final particle, so the total number of states in this
range is
\begin{equation}
 d\mathcal{N}_\beta=\left[V/(2\pi)^3\right]^{N_\beta}d\beta.
 \tag{3.4.8}\label{eq:3.4.8}
\end{equation}
Hence the total probability for the system to wind up in a range $d\beta$ of final states is
\begin{equation}
 dP(\alpha\to\beta)=P(\alpha\to\beta)d\mathcal{N}_\beta
 =\left[(2\pi)^3/V\right]^{N_\alpha}|S_{\beta\alpha}|^2d\beta.
 \tag{3.4.9}\label{eq:3.4.9}
\end{equation}

We will restrict our attention throughout this section to final states $\beta$ that are not only
different (however slightly) from the initial state $\alpha$, but that also satisfy the more
stringent condition, that no subset of the particles in the state $\beta$ (other than the whole
state itself) has precisely the same four-momentum as some corresponding subset of the particles
in the state $\alpha$. (In the language to be introduced in the next chapter, this means that
we are considering only the connected part of the $S$-matrix.) For such states, we may define
a delta function-free matrix element $M_{\beta\alpha}$:
\begin{equation}
 S_{\beta\alpha}=-2\pi i\,\delta_V^3(\mathbf{p}_\beta-\mathbf{p}_\alpha)
 \delta_T(E_\beta-E_\alpha)M_{\beta\alpha}.
 \tag{3.4.10}\label{eq:3.4.10}
\end{equation}
Our introduction of the box allows us to interpret the squares of delta functions in
$|S_{\beta\alpha}|^2$ for $\beta\ne\alpha$ as
\begin{align*}
 \left[\delta_V^3(\mathbf{p}_\beta-\mathbf{p}_\alpha)\right]^2
 &=\delta_V^3(\mathbf{p}_\beta-\mathbf{p}_\alpha)\delta_V^3(0)
 =\delta_V^3(\mathbf{p}_\beta-\mathbf{p}_\alpha)V/(2\pi)^3,\\
 \left[\delta_T(E_\beta-E_\alpha)\right]^2
 &=\delta_T(E_\beta-E_\alpha)\delta_T(0)
 =\delta_T(E_\beta-E_\alpha)T/2\pi,
\end{align*}
so Eq. \eqref{eq:3.4.9} gives a differential transition probability
\begin{align*}
 dP(\alpha\to\beta)={}&(2\pi)^2\left[(2\pi)^3/V\right]^{N_\alpha-1}
 (T/2\pi)|M_{\beta\alpha}|^2\\
 &\cdot\delta_V^3(\mathbf{p}_\beta-\mathbf{p}_\alpha)
 \delta_T(E_\beta-E_\alpha)d\beta.
\end{align*}
If we let $V$ and $T$ be very large, the delta function product here may be interpreted as an
ordinary four-dimensional delta function $\delta^4(p_\beta-p_\alpha)$. In this limit, the
transition probability is simply proportional to the time $T$ during which the interaction is
acting, with a coefficient that may be interpreted as a differential transition \emph{rate}:
\begin{align}
 d\Gamma(\alpha\to\beta)&\equiv dP(\alpha\to\beta)/T\notag\\
 &=(2\pi)^{3N_\alpha-2}V^{1-N_\alpha}|M_{\beta\alpha}|^2
 \delta^4(p_\beta-p_\alpha)d\beta,
 \tag{3.4.11}\label{eq:3.4.11}
\end{align}
where now
\begin{equation}
 S_{\beta\alpha}=-2\pi i\delta^4(p_\beta-p_\alpha)M_{\beta\alpha}.
 \tag{3.4.12}\label{eq:3.4.12}
\end{equation}
This is the master formula which is used to interpret calculations of $S$-matrix elements in
terms of predictions for actual experiments. We will come back to the interpretation of the
factor $\delta^4(p_\alpha-p_\beta)d\beta$ later in this section.

There are two cases of special importance:

\noindent$\boldsymbol{N_\alpha=1:}$

Here the volume $V$ cancels in Eq. \eqref{eq:3.4.11}, which gives the transition rate for a single-particle
state $\alpha$ to decay into a general multi-particle state $\beta$ as
\begin{equation}
 d\Gamma(\alpha\to\beta)=2\pi|M_{\beta\alpha}|^2
 \delta^4(p_\beta-p_\alpha)d\beta.
 \tag{3.4.13}\label{eq:3.4.13}
\end{equation}
Of course, this makes sense only if the time $T$ during which the interaction acts is much less
than the mean lifetime $\tau_\alpha$ of the particle $\alpha$, so we cannot pass to the limit
$T\to\infty$ in $\delta_T(E_\alpha-E_\beta)$. There is an unremovable width
$\Delta E\simeq1/T\gtrsim1/\tau_\alpha$ in this delta function, so Eq. \eqref{eq:3.4.13} is only useful
if the total decay rate $1/\tau_\alpha$ is much less than any of the characteristic energies of
the process.

\noindent$\boldsymbol{N_\alpha=2:}$

Here the rate \eqref{eq:3.4.11} is proportional to $1/V$, or in other words, to the density of either
particle at the position of the other one. Experimentalists generally report not the transition
rate per density, but the \emph{rate per flux}, also known as the \emph{cross-section}. The flux
of either particle at the position of the other particle is defined as the product of the density
$1/V$ and the relative velocity $u_\alpha$:
\begin{equation}
 \Phi_\alpha=u_\alpha/V.
 \tag{3.4.14}\label{eq:3.4.14}
\end{equation}
(A general definition of $u_\alpha$ is given below; for the moment we will content ourselves with
specifying that if either particle is at rest then $u_\alpha$ is defined as the velocity of the
other.) Thus the differential cross-section is
\begin{equation}
 d\sigma(\alpha\to\beta)\equiv d\Gamma(\alpha\to\beta)/\Phi_\alpha
 =(2\pi)^4u_\alpha^{-1}|M_{\beta\alpha}|^2
 \delta^4(p_\beta-p_\alpha)d\beta.
 \tag{3.4.15}\label{eq:3.4.15}
\end{equation}

Even though the cases $N_\alpha=1$ and $N_\alpha=2$ are the most important, transition rates
for $N_\alpha\geq3$ are all measurable in principle, and some of them are very important in
chemistry, astrophysics, etc. (For instance, in one of the main reactions that release energy
in the sun, two protons and an electron turn into a deuteron and a neutrino.) Section 3.6 presents
an application of the master transition rate formula \eqref{eq:3.4.11} for general numbers $N_\alpha$ of
initial particles.

We next take up the question of the Lorentz transformation properties of rates and cross-sections,
which will help us to give a more general definition of the relative velocity $u_\alpha$ in
Eq. \eqref{eq:3.4.15}. The Lorentz transformation rule \eqref{eq:3.3.1} for the $S$-matrix is complicated by the
momentum-dependent matrices associated with each particle's spin. To avoid this complication,
consider the absolute-value squared of \eqref{eq:3.3.1} (after factoring out the Lorentz-invariant delta
function in Eq. \eqref{eq:3.4.12}), and sum over all spins. The unitarity of the matrices $D^{(j)}(W)$
(or their analogs for zero mass) then shows that, apart from the energy factors in \eqref{eq:3.3.1}, the
sum is Lorentz-invariant. That is, the quantity
\begin{equation}
 \sum_{\mathrm{spins}}|M_{\beta\alpha}|^2\prod_\beta E\prod_\alpha E
 \equiv R_{\beta\alpha}
 \tag{3.4.16}\label{eq:3.4.16}
\end{equation}
is a scalar function of the four-momenta of the particles in states $\alpha$ and $\beta$.
(By $\prod_\alpha E$ and $\prod_\beta E$ is meant the product of all the single-particle energies
$p^0=\sqrt{\mathbf{p}^2+m^2}$ for the particles in the states $\alpha$ and $\beta$.)

We can now write the spin-summed single-particle decay rate \eqref{eq:3.4.13} as
\[
 \sum_{\mathrm{spins}}d\Gamma(\alpha\to\beta)
 =2\pi E_\alpha^{-1}R_{\beta\alpha}\delta^4(p_\beta-p_\alpha)d\beta
 /\prod_\beta E.
\]
The factor $d\beta/\prod_\beta E$ may be recognized as the product of the Lorentz-invariant
momentum-space volume elements \eqref{eq:2.5.15}, so it is Lorentz-invariant. So also are
$R_{\beta\alpha}$ and $\delta^4(p_\beta-p_\alpha)$, leaving just the non-invariant factor
$1/E_\alpha$, where $E_\alpha$ is the energy of the single initial particle. Our conclusion
then is that the decay rate has the same Lorentz transformation property as $1/E_\alpha$.
This is, of course, just the usual special-relativistic time dilation---the faster the particle,
the slower it decays.

Similarly, our result \eqref{eq:3.4.15} for the spin-summed cross-section may be written as
\[
 \sum_{\mathrm{spins}}d\sigma(\alpha\to\beta)
 =(2\pi)^4u_\alpha^{-1}E_1^{-1}E_2^{-1}R_{\beta\alpha}
 \delta^4(p_\alpha-p_\beta)d\beta/\prod_\beta E,
\]
where $E_1$ and $E_2$ are the energies of the two particles in the initial state $\alpha$.
It is conventional to define the cross-section to be (when summed over spins) a Lorentz-invariant
function of four-momenta. The factors $R_{\beta\alpha}$, $\delta^4(p_\beta-p_\alpha)$, and
$d\beta/\prod_\beta E$ are already Lorentz-invariant, so this means that we must define the
relative velocity $u_\alpha$ in arbitrary inertial frames so that $u_\alpha E_1E_2$ is a scalar.
We also mentioned earlier that in the Lorentz frame in which one particle (say, particle 1) is
at rest, $u_\alpha$ is the velocity of the other particle. This uniquely determines $u_\alpha$
to have the value in general Lorentz frames
\begin{equation}
 u_\alpha=\sqrt{(p_1\cdot p_2)^2-m_1^2m_2^2}\big/E_1E_2,
 \tag{3.4.17}\label{eq:3.4.17}
\end{equation}
where $p_1,p_2$ and $m_1,m_2$ are the four-momenta and mass of the two particles in the initial
state $\alpha$.

As a bonus, we note that in the `center-of-mass' frame, where the total three-momentum vanishes,
we have
\[
 p_1=(E_1,\mathbf{p}),\qquad p_2=(E_2,-\mathbf{p}),
\]
and here Eq. \eqref{eq:3.4.17} gives
\begin{equation}
 u_\alpha=\frac{|\mathbf{p}|(E_1+E_2)}{E_1E_2}
 =\left|\frac{\mathbf{p}_1}{E_1}-\frac{\mathbf{p}_2}{E_2}\right|
 \tag{3.4.18}\label{eq:3.4.18}
\end{equation}
as might have been expected for a relative velocity. However, in this frame $u_\alpha$ is not
really a physical velocity; in particular, Eq. \eqref{eq:3.4.18} shows that for extremely relativistic
particles, it can take values as large as 2.
\footnote{Eq. \eqref{eq:3.4.17} makes it obvious that $E_1E_2u_\alpha$ is a scalar. Also, when particle 1
is at rest, we have $\mathbf{p}_1=0$, $E_1=m_1$, so $p_1\cdot p_2=-m_1E_2$, and so Eq. \eqref{eq:3.4.17}
gives $u_\alpha=\sqrt{E_2^2-m_2^2}/E_2=|\mathbf{p}_2|/E_2$, which is just the velocity of
particle 2.}

We now turn to the interpretation of the so-called phase-space factor
$\delta^4(p_\beta-p_\alpha)d\beta$, which appears in the general formula \eqref{eq:3.4.11} for transition
rates, and also in Eqs. \eqref{eq:3.4.13} and \eqref{eq:3.4.15} for decay rates and cross-sections. We here
specialize to the case of the `center-of-mass' Lorentz frame, where the total three-momentum
of the initial state vanishes
\begin{equation}
 \mathbf{p}_\alpha=0.
 \tag{3.4.19}\label{eq:3.4.19}
\end{equation}
(For $N_\alpha=1$, this is just the case of a particle decaying at rest.) If the final state
consists of particles with momenta $\mathbf{p}'_1,\mathbf{p}'_2,\ldots$, then
\begin{equation}
 \delta^4(p_\beta-p_\alpha)d\beta
 =\delta^3(\mathbf{p}'_1+\mathbf{p}'_2+\cdots)
 \delta(E'_1+E'_2+\cdots-E)d^3p'_1d^3p'_2\cdots,
 \tag{3.4.20}\label{eq:3.4.20}
\end{equation}
where $E\equiv E_\alpha$ is the total energy of the initial state. Any one of the
$\mathbf{p}'_k$ integrals, say over $\mathbf{p}'_1$, can be done trivially by just dropping the
momentum delta function
\begin{equation}
 \delta^4(p_\beta-p_\alpha)d\beta\longrightarrow
 \delta(E'_1+E'_2+\cdots-E)d^3p'_2\cdots
 \tag{3.4.21}\label{eq:3.4.21}
\end{equation}
with the understanding that wherever $\mathbf{p}'_1$ appears (as in $E'_1$) it must be replaced
with
\begin{equation}
 \mathbf{p}'_1=-\mathbf{p}'_2-\mathbf{p}'_3-\cdots.
 \tag{3.4.22}\label{eq:3.4.22}
\end{equation}
We can similarly use the remaining delta function to eliminate any one of the remaining integrals.

In the simplest case, there are just two particles in the final state. Here \eqref{eq:3.4.21} gives
\[
 \delta^4(p_\beta-p_\alpha)d\beta\longrightarrow
 \delta(E'_1+E'_2-E)d^3p'_2.
\]
In more detail, this is
\begin{equation}
 \delta^4(p_\beta-p_\alpha)d\beta\longrightarrow
 \delta\left(\sqrt{|\mathbf{p}'_1|^2+m_1'^2}
 +\sqrt{|\mathbf{p}'_1|^2+m_2'^2}-E\right)
 |\mathbf{p}'_1|^2d|\mathbf{p}'_1|d\Omega,
 \tag{3.4.23}\label{eq:3.4.23}
\end{equation}
where $\mathbf{p}'_2=-\mathbf{p}'_1$ and $d\Omega\equiv\sin\theta\,d\theta\,d\phi$ is the
solid-angle differential for $\mathbf{p}'_1$. This can be simplified by using the standard formula
\[
 \delta(f(x))=\delta(x-x_0)/|f'(x_0)|,
\]
where $f(x)$ is an arbitrary real function with a single simple zero at $x=x_0$. In our case,
the argument $E'_1+E'_2-E$ of the delta function in Eq. \eqref{eq:3.4.23} has a unique zero at
$|\mathbf{p}'_1|=k'$, where
\begin{equation}
 k'=\frac{\sqrt{(E^2-m_1'^2-m_2'^2)^2-4m_1'^2m_2'^2}}{2E},
 \tag{3.4.24}\label{eq:3.4.24}
\end{equation}
\begin{equation}
 E'_1=\sqrt{k'^2+m_1'^2}=\frac{E^2-m_2'^2+m_1'^2}{2E},
 \tag{3.4.25}\label{eq:3.4.25}
\end{equation}
\begin{equation}
 E'_2=\sqrt{k'^2+m_2'^2}=\frac{E^2-m_1'^2+m_2'^2}{2E},
 \tag{3.4.26}\label{eq:3.4.26}
\end{equation}
with derivative
\begin{align}
 \left[\frac{d}{d|\mathbf{p}'_1|}
 \left(\sqrt{|\mathbf{p}'_1|^2+m_1'^2}
 +\sqrt{|\mathbf{p}'_1|^2+m_2'^2}-E\right)\right]_{|\mathbf{p}'_1|=k'}
 &=\frac{k'}{E'_1}+\frac{k'}{E'_2}=\frac{k'E}{E'_1E'_2}.
 \tag{3.4.27}\label{eq:3.4.27}
\end{align}
We can thus drop the delta function and the differential $d|\mathbf{p}'_1|$ in Eq. \eqref{eq:3.4.23},
by dividing by \eqref{eq:3.4.27},
\begin{equation}
 \delta^4(p_\beta-p_\alpha)d\beta\longrightarrow
 \frac{k'E'_1E'_2}{E}\,d\Omega
 \tag{3.4.28}\label{eq:3.4.28}
\end{equation}
with the understanding that $k',E'_1$, and $E'_2$ are given everywhere by Eqs. \eqref{eq:3.4.24}--\eqref{eq:3.4.26}.
In particular, the differential rate \eqref{eq:3.4.13} for decay of a one-particle state of zero momentum
and energy $E$ into two particles is
\begin{equation}
 \frac{d\Gamma(\alpha\to\beta)}{d\Omega}
 =\frac{2\pi k'E'_1E'_2}{E}|M_{\beta\alpha}|^2
 \tag{3.4.29}\label{eq:3.4.29}
\end{equation}
and the differential cross-section for the two-body scattering process $12\to1'2'$ is given by
Eq. \eqref{eq:3.4.15} as
\begin{equation}
 \frac{d\sigma(\alpha\to\beta)}{d\Omega}
 =\frac{(2\pi)^4k'E'_1E'_2}{Eu_\alpha}|M_{\beta\alpha}|^2
 =\frac{(2\pi)^4k'E'_1E'_2E_1E_2}{E^2k}|M_{\beta\alpha}|^2,
 \tag{3.4.30}\label{eq:3.4.30}
\end{equation}
where $k\equiv|\mathbf{p}_1|=|\mathbf{p}_2|$.

The above case $N_\beta=2$ is particularly simple, but there is one nice result for $N_\beta=3$
that is also worth recording. For $N_\beta=3$, Eq. \eqref{eq:3.4.21} gives
\[
 \delta^4(p_\beta-p_\alpha)d\beta\longrightarrow d^3p'_2d^3p'_3
 \cdot\delta\left(\sqrt{|\mathbf{p}'_2+\mathbf{p}'_3|^2+m_1'^2}
 +\sqrt{|\mathbf{p}'_2|^2+m_2'^2}+\sqrt{|\mathbf{p}'_3|^2+m_3'^2}-E\right).
\]
We write the momentum-space volume as
\[
 d^3p'_2d^3p'_3=|\mathbf{p}'_2|^2d|\mathbf{p}'_2|
 |\mathbf{p}'_3|^2d|\mathbf{p}'_3|d\Omega_3d\phi_{23}d\cos\theta_{23},
\]
